\scanpage{145}
Nous allons étudier (14) pour $i=0,1$.

1) $H^0_{\mathrm{DR}}(S,E)$ est formée des sections \inserted{$x$} de $E$ qui sont
``horizontales pour la connexion'' $\delta$, ce qui \deleted{s'écrit par}\inserted{signifie par}
définition \deleted{$d_E^0x=0$}
\[
d_E^0(x)=0\qquad\text{i.e.}\qquad\delta_E x=d_E^1x \tag*{(16)}
\]
Alors (14) est l'injection canonique.

% slanted
\margin{Comparer Cat. \unsure{cof.|cat.|cf.} \unsure{additives|abéliennes|additions} 1.8.}

2) Soit $K^{\unsure{\cdot|\bullet}}:0\longrightarrow K^0\xrightarrow{d^0}K^1\xrightarrow{d^1}K^2\longrightarrow\cdots$
un complexe \deleted{\unsure{On}}\inserted{sur} \deleted{de faisceaux ab.}
\deleted{constate que $\mathbb H^1(S$}
de faisceaux abéliens sur un topos \inserted{$S$}. On constate que
$\mathbb H^1(S,K^{\cdot})$ s'interprète comme le groupe des
classes d'isomorphie \deleted{de} d'objets de la catégorie suivante \inserted{$C(K^{\cdot})$} :

$\operatorname{Ob} C(K^{\cdot})=$ triples $(P^0,\zeta)$, avec $P$ un torseur sous $K^0$,
$\zeta$ une section de $P^1=P\mathbin{\overset{K^0}{\times}}K^1$, telle que la section
\deleted{\unsure{canoni}} $d^1_*(\zeta)$ de $P^2=P\mathbin{\overset{K^0}{\times}}K^2\simeq
P^1\mathbin{\overset{K^1}{\times}}K^2$ soit la section
canonique, \deleted{\unsure{déi}} déduit du fait que $d^1d^0=0$.

$\mathrm{Fl}\ C(K^{\cdot})$ : on devine

\deleted{L'hom canonique $\mathbb H^1(S,K^{\cdot})\longrightarrow H^1(\unsure{S,}$}\inserted{Cette description de $\mathbb H^1(S,K^{\cdot})$ est compatible avec}
les hom. \unsure{en|entre} les complexes $K^{\cdot}$. Par suite,
l'hom \inserted{can.} \deleted{$\mathbb H^1(S,K^{\cdot}$}
\[
\mathbb H^1(S,K^{\cdot})\longrightarrow H^1(S,K^0)
\]
déduit de l'\deleted{\unsure{augmentation de $K$}}\inserted{hom troncature} $K^{\cdot}\longrightarrow K^0$ correspond au
foncteur $(P^0,\zeta)\longmapsto P^0$ ``oubli de $\zeta$''.

Noter que \deleted{la donnée, pour}\inserted{pour un} $P^0$ donné (torseur sous $K^0$) la
donnée d'une trivialisation $\zeta$ de $P^1$ équiv. à un isom
$P^1\simeq K^1$, donc les $\zeta$ correspondent \deleted{aux hom}
% archive sheet number: 144
